\begin{papercard}{oconnor2020}
\cf{Stages} \stagetag{C} \stagetag{I}. A practical review of whether and how to assess \emph{intercoder reliability} (ICR) in qualitative analysis. It covers arguments for and against, a decision list (coders, share of data, independence, segmentation, number and depth of codes, statistic, reporting, interpretation, follow-up) and a suggested procedure. This is the main source for the thesis' two-coder design.

\cf{Definitions} ICR is a numerical measure of agreement between coders on how the same data should be coded. Strictly, \emph{interrater} reliability applies to ordinal or interval ratings, whereas ICR applies to nominal categories. \emph{Intracoder} reliability means one coder's consistency over time. \emph{Intercoder consistency} is unquantified comparison and discussion \pc{oconnor2020}{2}{ICR is a numerical measure of the agreement between different coders}, \pc{oconnor2020}{2}{However, technically IRR refers to cases where data are rated on some ordinal or interval scale}. ICR belongs to the coding phase, because that is where analysis first leaves the raw data. Theme construction afterwards may legitimately differ between analysts if the rationale is transparent \pc{oconnor2020}{2}{The logic of applying ICR to the earlier coding phase is that coding is the first place}.

\cf{Current practice} Percentage agreement was the most common method yet is inadequate; reporting is often vague, with 69\,\% of content-analysis articles mentioning ICR and many using inappropriate methods \pc{oconnor2020}{3}{Lombard et al. (2002) report that of content analysis articles published in the mass communications}. The field is moving to Krippendorff's $\alpha$ and similar statistics.

\cf{Arguments for} (1) Rigour and transparency of the coding frame, i.e.\ communicability across persons \pc{oconnor2020}{3}{The most commonly cited rationale for performing an ICR assessment is to assess the rigor}. (2) An external ``badge of trustworthiness'' \pc{oconnor2020}{3}{ICR can thus serve as a badge of trustworthiness}. (3) Internal benefits: consistent decisions across multiple coders (also cross-language), self-discipline, and \emph{reflexivity and dialogue}, since the content of disagreements can be more valuable than the final score. Example: codes such as ``self-control'' proved unclear to a second coder and were redefined \pc{oconnor2020}{4}{the content of inter- coder disagreements can be equally, if not more valuable than the ultimate degree of consistency}. (4) Consequences: findings that inform policy or clinical practice need confidence in the evidence base.

\cf{Objections and replies} ICR may contradict interpretive epistemology (Braun \& Clarke, Stenbacka, Madill) \pc{oconnor2020}{4}{Perhaps the most frequently aired objection is that ICR essentially contradicts the interpretative}. Reply: agreement shows that a group within a common framework can reach consensual interpretation, not that a single true meaning exists \pc{oconnor2020}{5}{Rather, it shows that a group of researchers working within a common conceptual framework can reach a consensual}. Risk of \emph{false precision}: nonsensical maths can inflate perceived quality, and students may fixate on ICR \pc{oconnor2020}{5}{Research shows inclusion of entirely nonsensical mathematical information can inflate judgments}. ICR is ``never an end in itself''.

\cf{Practical decisions} \textbf{Manual or electronic}: built-in ICR tools in NVivo, Dedoose and QDA Miner are often opaque and sensitive to trivial differences, so exporting to statistics software gives control \pc{oconnor2020}{5}{However, in the authors’ experience, the process by which automated ICR tools calculate an ICR figure}. \textbf{How many coders}: at least two independent coders; a new coder needs a new check; when one person codes the rest, a second coder codes a sample \pc{oconnor2020}{5}{A minimum of two independent coders is necessary to establish ICR}. \textbf{Share of data}: typically 10--25\,\% of units, drawn at random or stratified, after a small informal pre-check (one interview) \pc{oconnor2020}{5}{Depending on the size of the data set, 10–25\% of data units would be typical}. Repeated rounds are possible but may only reflect ``interpretive convergence'' \pc{oconnor2020}{6}{However, it is possible this reliability may simply reflect “interpretive convergence”}. \textbf{Independence}: independent double-coding without conferral; discussion before coding is a design choice, and external coders raise data-protection issues \pc{oconnor2020}{6}{It is generally accepted that the physical double-coding should be performed independently without conferral}. Training is needed more for \emph{latent} than for surface codes.

\cf{Segmentation} Whole document, prespecified units (paragraph, response, sentence) or ad hoc conceptual segments. Larger units preserve context (validity) but lower reliability and sensitivity \pc{oconnor2020}{6}{In relation to ICR specifically, longer text units are usually associated with poorer reliability}. Ad hoc units create the ``unitization problem'': ``lumpers'' vs.\ ``splitters''. Solution (after Campbell et al.): the primary coder segments and codes, then a \emph{clean file with segments but without codes} goes to the second coder \pc{oconnor2020}{7}{Once the first coder has saved the coded transcript, they can then create a second document where the data segments are visible}.

\cf{Code number and depth} More codes mean lower ICR. Limits of 30--40 (MacQueen) or about 20 (Hruschka); alternatively group codes into families and take family-specific passes \pc{oconnor2020}{7}{MacQueen et al. (1998) suggest researchers concerned with achieving satisfactory reliability should work with an upper limit of 30–40 codes}. Never trade away comprehensiveness or theory for reliability. Define explicitly, add qualifications and exclusions, and assume nothing is obvious \pc{oconnor2020}{8}{Theoretical relevance or meaning should never be sacrificed for reliability}.

\cf{Statistic and reporting} Percentage agreement is rejected because it ignores chance and is ill-defined for multi-coding. Use Cohen's $\kappa$, Krippendorff's $\alpha$, Scott's $\pi$, Fleiss' K etc. Cohen's $\kappa$ needs exclusive coding, and $\alpha$ is flexible with more than two coders \pc{oconnor2020}{8}{The pri- mary advantage these statistics offer over percentage agreement is correction for the probability}. Chi-square, Cronbach's $\alpha$ and Pearson $r$ are inappropriate. Export each coder's 0/1 code matrix to SPSS or R, one row per unit and one column per code per coder (Figure 2) \pc{oconnor2020}{9}{If a code has been applied to a data unit, the relevant cell shows 1}. Report per code, not only pooled, because a pooled value hides poor codes; at least give the range and distribution \pc{oconnor2020}{9}{Second, pooling all codes’ reliability figures means codes with poor reliability can be “hidden”}.

\cf{Interpretation and follow-up} Statistical significance is never evidence of ICR. Rules of thumb are arbitrary (Landis \& Koch; Neuendorf: $>.9$ acceptable to all, $>.8$ to many) and \emph{must be stipulated in advance}, with stricter thresholds for high-stakes work \pc{oconnor2020}{9}{Whatever interpretative framework is chosen should be stipulated in advance and not decided post hoc}. Follow-up options for weak codes: discard, revise and re-test, or keep theoretically important codes. For final data, use majority rule, the senior coder's call, or \emph{negotiated agreement} (Campbell: 54\,\% to 96\,\%) \pc{oconnor2020}{9}{Campbell et al. (2013) describe such a strategy of “negotiated agreement”}.

\cf{Suggested procedure (Figure 3)} A priori: number of coders, amount of duplicate data, unit, depth, statistic and threshold. Then familiarise, draft the frame, Coder 1 segments and codes the ICR subset, a clean file goes to Coder 2, informal check, export and compute per code, compare with the threshold, revise on a \emph{different} subset, then code the rest \pc{oconnor2020}{10}{The revised coding frame can be evaluated using the same process, preferably on a different subset of data}.

\cf{Use for this thesis} A near-complete protocol for the two-coder stage. (1) Fix coders, subset fraction (e.g.\ about 10--25\,\% of walkthrough steps or apps), unit of analysis, statistic and threshold \emph{before} piloting and record them in the Decision Log. (2) Because the walkthrough steps are pre-segmented by the protocol, the unitization problem largely disappears, which is worth stating. (3) Use the clean-file hand-off, and compute $\kappa$ or $\alpha$ \emph{per code} with a 0/1 matrix. (4) Report disagreement handling (negotiated agreement) and re-test on a fresh subset after revisions. (5) Keep the sceptical position in view: say what ICR does and does not show (coding-frame communicability, not truth). It does not settle which statistic to use under prevalence skew (see Gwet, Feinstein \& Cicchetti).
\end{papercard}
